\section{Results}
\label{sec:results}
\label{sec:5}

\subsection{Expected parallax over India}\label{sec:5.1}

The parallax shift grows with distance from the sub-satellite point at 74°E (Fig.~\ref{fig:1}). For a 12 km cloud top it is about 2.5--3 km over the southern peninsula (VZA 12--15°), 5--7 km over central India (VZA 22--30°) and 8--9 km over the Himalaya and Northeast India (VZA 33--38°); for a 16 km top it reaches 11--12 km there. The shift points away from the satellite, which over India is roughly northwards; over the Northeast it turns towards the north-east (VAA up to 217°).

In the validation cases the correction moved deep convective cores (BT \textless{} 220 K) by a median of 4.2 km in South India, 7.3 km in Central India, 9.8 km over the Himalaya and 10.5 km over Northeast India (Table~\ref{tab:2}). These shifts follow tan(VZA) as expected; the solver itself agrees with pyproj, pyorbital and satpy to within 1 mm, 0.01° and 25 m respectively (Section~\ref{sec:3}).

\subsection{INSAT-3DR navigation}\label{sec:5.2}

Coastline matching shows that INSAT-3DR L1C positions carry a navigation error that changes with the season (Fig.~\ref{fig:2}). From late September to mid-March, positions need a correction of 7--20 km to the west; from mid-April to September the error is small. The switch near the equinoxes is consistent with the satellite's twice-yearly yaw flip.

During the monsoon the error cannot be measured reliably for a single day. Night-time matching fails because land and sea have almost the same TIR1 temperature (r \textless{} 0.4 for every night-time case scan). Daytime matching works, but the same coast in scans 1--2 h apart differs by about 9--10 km, and the best correlation is typically only 0.01--0.05 above the unshifted one. Of 199 daytime matches near the validation storms, 73 have r ≥ 0.6 and lie inside the search range. Pooling the 36 matches with r ≥ 0.7 gives a monsoon correction of −0.5 km east and +2.5 km north, which we adopt, with an uncertainty of about ±2 km. The median over all monsoon measurements across India (+1.8 km east, +4.8 km north) is used as a sensitivity test.

\subsection{Validation cases}\label{sec:5.3}

Of 33 coincident INSAT--GPM cases (GPM within 5 min of the INSAT scan at the storm; for the usable cases, median 1.3 min and at most 2.5 min), 30 could be processed. Of these, 9 had too little convection (fewer than 100 convective footprints in the main cluster) and 7 gave no reliable match; 14 cases were usable (Table~\ref{tab:1}). For the field measure, 12 cases give a reliable match in every correction state, and these 12 are used throughout.

\begin{longtable}[]{@{}
  >{\raggedright\arraybackslash}p{(\linewidth - 14\tabcolsep) * \real{0.1250}}
  >{\raggedright\arraybackslash}p{(\linewidth - 14\tabcolsep) * \real{0.1250}}
  >{\raggedleft\arraybackslash}p{(\linewidth - 14\tabcolsep) * \real{0.1250}}
  >{\raggedleft\arraybackslash}p{(\linewidth - 14\tabcolsep) * \real{0.1250}}
  >{\raggedleft\arraybackslash}p{(\linewidth - 14\tabcolsep) * \real{0.1250}}
  >{\raggedleft\arraybackslash}p{(\linewidth - 14\tabcolsep) * \real{0.1250}}
  >{\raggedleft\arraybackslash}p{(\linewidth - 14\tabcolsep) * \real{0.1250}}
  >{\raggedleft\arraybackslash}p{(\linewidth - 14\tabcolsep) * \real{0.1250}}@{}}
\caption{Validation cases.}\label{tab:1}\\
\toprule\noalign{}
\begin{minipage}[b]{\linewidth}\raggedright
Region
\end{minipage} & \begin{minipage}[b]{\linewidth}\raggedright
VZA at storm
\end{minipage} & \begin{minipage}[b]{\linewidth}\raggedleft
Cases
\end{minipage} & \begin{minipage}[b]{\linewidth}\raggedleft
Not run (missing files)
\end{minipage} & \begin{minipage}[b]{\linewidth}\raggedleft
Processed
\end{minipage} & \begin{minipage}[b]{\linewidth}\raggedleft
Usable
\end{minipage} & \begin{minipage}[b]{\linewidth}\raggedleft
Too little convection
\end{minipage} & \begin{minipage}[b]{\linewidth}\raggedleft
No reliable match
\end{minipage} \\*
\midrule\noalign{}
\endfirsthead
\bottomrule\noalign{}
\endlastfoot
Northeast India & 31--39° & 11 & 1 & 10 & 5 & 3 & 2 \\*
Himalaya & 29--37° & 12 & 1 & 11 & 4 & 6 & 1 \\*
Central India & 24--26° & 4 & 0 & 4 & 2 & 0 & 2 \\*
South India & 14--21° & 6 & 1 & 5 & 3 & 0 & 2 \\*
\textbf{Total} & & \textbf{33} & \textbf{3} & \textbf{30} & \textbf{14} & \textbf{9} & \textbf{7} \\*
\end{longtable}

\subsection{Effect of parallax correction}\label{sec:5.4}

Parallax correction brings INSAT closer to GPM on average (Fig.~\ref{fig:3}, Table~\ref{tab:2}). For the storm-top field within 250 km of the main convective cluster, the median absolute offset falls from 5.5 km to 3.1 km and 10 of the 12 cases improve. The mean offset changes from +4.3 km to −0.9 km. Case by case, the absolute offset falls by a median of 1.6 km (95 \% bootstrap interval 0.6--5.4 km); the improvement is significant by the sign test (p = 0.039) but not by the Wilcoxon signed-rank test (p = 0.13), so the evidence is moderate.

The correction acts only along the line of sight, as it should: the cross-look offset is the same before and after (median absolute 8.8 and 9.0 km; median change 0.0 km, 95 \% interval −0.3 to +0.2 km). The cross-look scatter is therefore not a parallax effect but reflects storm structure, storm motion within the 1--2 min between the two observations, and the different quantities the two instruments see.

The gain is clearest at high viewing angles (Table~\ref{tab:2}). Over the Himalaya and Northeast India, where the deep cores moved by about 10 km, the mean offset moves from +2.0 and +5.2 km to −4.9 and −1.4 km. In South India, where the parallax is only 3--5 km, it changes from +6.2 to +4.9 km, within the case-to-case scatter.

\begin{longtable}[]{@{}
  >{\raggedright\arraybackslash}p{(\linewidth - 12\tabcolsep) * \real{0.1429}}
  >{\raggedleft\arraybackslash}p{(\linewidth - 12\tabcolsep) * \real{0.1429}}
  >{\raggedleft\arraybackslash}p{(\linewidth - 12\tabcolsep) * \real{0.1429}}
  >{\raggedleft\arraybackslash}p{(\linewidth - 12\tabcolsep) * \real{0.1429}}
  >{\raggedleft\arraybackslash}p{(\linewidth - 12\tabcolsep) * \real{0.1429}}
  >{\raggedleft\arraybackslash}p{(\linewidth - 12\tabcolsep) * \real{0.1429}}
  >{\raggedleft\arraybackslash}p{(\linewidth - 12\tabcolsep) * \real{0.1429}}@{}}
\caption{Line-of-sight offset (storm-top field, 250 km window) by region; mean over the cases valid in every state.}\label{tab:2}\\
\toprule\noalign{}
\begin{minipage}[b]{\linewidth}\raggedright
Region
\end{minipage} & \begin{minipage}[b]{\linewidth}\raggedleft
Cases
\end{minipage} & \begin{minipage}[b]{\linewidth}\raggedleft
Mean VZA
\end{minipage} & \begin{minipage}[b]{\linewidth}\raggedleft
Applied shift at deep cores, median {[}km{]}
\end{minipage} & \begin{minipage}[b]{\linewidth}\raggedleft
Raw {[}km{]}
\end{minipage} & \begin{minipage}[b]{\linewidth}\raggedleft
Parallax {[}km{]}
\end{minipage} & \begin{minipage}[b]{\linewidth}\raggedleft
Navigation + parallax {[}km{]}
\end{minipage} \\*
\midrule\noalign{}
\endfirsthead
\bottomrule\noalign{}
\endlastfoot
South India & 3 & 15° & 4.2 & +6.2 & +4.9 & +7.2 \\*
Central India & 1 & 24° & 7.3 & +3.5 & −0.5 & +1.9 \\*
Himalaya & 4 & 34° & 9.8 & +2.0 & −4.9 & −2.6 \\*
Northeast India & 4 & 34° & 10.5 & +5.2 & −1.4 & −0.2 \\*
\textbf{All} & \textbf{12} & & 9.8 & \textbf{+4.3} & \textbf{−0.9} & \textbf{+1.0} \\*
\end{longtable}

Applied shifts are medians over all 14 usable cases; offsets use the 12 cases valid in every state.

\subsection{Parallax-only correction overshoots at storm cores}\label{sec:5.5}

If the uncorrected offset were pure parallax, it would grow with tan(VZA), by about 10--12 km per unit tan(VZA) for deep convection. Across the 12 cases no such growth is visible (Spearman ρ = −0.03), but with so few cases the slope is poorly constrained (−4 km per unit tan(VZA), 95 \% bootstrap interval −19 to +15 km), and a parallax-like trend cannot be excluded. The two offset measures give a clearer picture.

\textbf{Cell by cell} (Fig.~\ref{fig:4}), uncorrected INSAT cold cores already sit almost on top of GPM's deep-echo cells (median −0.1 km over 69 cells in 11 cases), and parallax correction moves them 5.7 km past (median −5.7 km; 95 \% interval −8.2 to −2.2 km from a bootstrap that resamples cases, since cells of one storm are not independent). The \textbf{field correlation} sees the opposite sign before correction (+4.3 km) and lands near zero after it (−0.9 km). The two measures differ by a steady 4--5 km in every state. The field measure is driven by the large cold anvil against the rain area, the cell measure by the coldest tops against the tallest echoes.

Three sources of this behaviour were tested:

\begin{enumerate}
\def\labelenumi{\arabic{enumi}.}
\tightlist
\item
  \textbf{GPM's slanted beam.} Placing storm tops on the slanted Ku beam instead of at the surface footprint (a median displacement of 1 km, up to 5 km) changes the cell medians by at most 0.6 km. It adds scatter, not bias.
\item
  \textbf{Cloud-top height.} In 95 matched cells, the BT-derived height of the INSAT cold core is 16.9 km (median), against a GPM maximum echo top of 12.45 km (Fig.~\ref{fig:5}); about half the cores sit at the tropopause. ERA5 profiles instead of the standard atmosphere give 16.75 km and change every offset by under 0.5 km. The gap between infrared cloud top and radar echo top is therefore physical: radar echo tops lie 1--3 km below the infrared top of deep convection. It accounts for about 1--2 km of the cell-measure overshoot.
\item
  \textbf{Navigation.} Adding the monsoon navigation correction moves the cell median from −5.7 to −4.1 km (per-case change +1.3 km, 95 \% interval +0.5 to +1.9 km, Wilcoxon p = 0.002) and the field mean from −0.9 to +1.0 km (Fig.~\ref{fig:4}). Navigation thus accounts for about 1.5--3 km of the overshoot, depending on which monsoon estimate is used; the remaining cell overshoot (−4.1 km, interval −7.8 to −0.6 km) is still significant.
\end{enumerate}

With navigation and parallax corrected together, the storm-top field gives the smallest median offset of any state: mean +1.0 km, median absolute 1.8 km (Table~\ref{tab:3}). This is a point estimate: with 12 cases it is not significantly better than parallax correction alone (4 of 12 cases improve; median change +1.1 km, 95 \% interval −1.9 to +1.4 km; Wilcoxon p = 0.91), and its interval is wide (1.4--8.6 km). The cell measure remains at −4.1 km, so the two measures bracket zero, as they do only in the corrected states.

\begin{longtable}[]{@{}
  >{\raggedright\arraybackslash}p{(\linewidth - 8\tabcolsep) * \real{0.2000}}
  >{\raggedleft\arraybackslash}p{(\linewidth - 8\tabcolsep) * \real{0.2000}}
  >{\raggedleft\arraybackslash}p{(\linewidth - 8\tabcolsep) * \real{0.2000}}
  >{\raggedleft\arraybackslash}p{(\linewidth - 8\tabcolsep) * \real{0.2000}}
  >{\raggedleft\arraybackslash}p{(\linewidth - 8\tabcolsep) * \real{0.2000}}@{}}
\caption{Offsets in each correction state (12 cases for the field measure; 69--75 cells for the cell measure).}\label{tab:3}\\
\toprule\noalign{}
\begin{minipage}[b]{\linewidth}\raggedright
State
\end{minipage} & \begin{minipage}[b]{\linewidth}\raggedleft
Field mean {[}km{]}
\end{minipage} & \begin{minipage}[b]{\linewidth}\raggedleft
Field median absolute {[}km{]}
\end{minipage} & \begin{minipage}[b]{\linewidth}\raggedleft
Cases improved
\end{minipage} & \begin{minipage}[b]{\linewidth}\raggedleft
Cells, median {[}km{]}
\end{minipage} \\*
\midrule\noalign{}
\endfirsthead
\bottomrule\noalign{}
\endlastfoot
Raw & +4.3 & 5.5 & -- & −0.1 \\*
Navigation only & +6.3 & 7.6 & -- & +0.1 \\*
Parallax only & −0.9 & 3.1 & 10 of 12 & −5.7 \\*
Navigation + parallax & \textbf{+1.0} & \textbf{1.8} & 7 of 12 & −4.1 \\*
All-India navigation + parallax (sensitivity) & +3.8 & 4.7 & 6 of 12 & −3.3 \\*
\end{longtable}

\subsection{Terrain}\label{sec:5.6}

Terrain-aware correction with SRTM moves clear-sky ground by up to 3--4 km over the Himalaya and the Tibetan Plateau, 1--2 km over the Northeast hills and under 0.5 km over the Western Ghats and the Deccan (Fig.~\ref{fig:6}). In the Himalayan validation scenes, clear-sky ground pixels move by 1.3--2.5 km (median) and up to 5.8 km, and 9--19 \% of grid cells change.

At the storms, terrain makes no difference: the field and cell offsets are identical to 0.1 km with and without it. Storm tops lie 10--17 km above the ground, so the height of the ground beneath them does not enter the correction.

\subsection{Benchmark}\label{sec:5.7}

The full pipeline processed all 148 available monsoon scenes without a failure. Per scene (reading two channels over 0--40°N, 60--100°E, cloud-top height, and the corrected and uncorrected grids at 0.04°) it took 1.8 s (median; 3.2 s at most) on a laptop, with about 450 MB of memory.

After correction, 98.7 \% of grid cells hold an observed pixel, 1.3 \% are filled from neighbours and 0.007 \% stay empty behind tall clouds. Corrected clouds cover 2.5 \% of previously clear cells (median). All other clear cells are identical in the corrected and uncorrected grids, so the two datasets differ only where clouds moved.

Fourteen scenes from 11--14 June 2023 turned out to be partial INSAT-3DR scans, with up to 57 \% of the image missing. The pipeline handles them; the two validation cases on those dates are unaffected because their storms lie inside the scanned part.

